\chapter{An Affair of State}

\lettrine{T}{he} cardinal, on passing into his cabinet, found the  Comte de la Fère, who was waiting for him, engaged in admiring a very fine  Raphael placed over a sideboard covered with a plate. His eminence came in  softly, lightly, and as silently as a shadow, and surprised the countenance of  the comte, as he was accustomed to do, pretending to divine by the simple  expression of the face of his interlocutor what would be the result of the  conversation.

But this time Mazarin was foiled in his expectation: he read nothing upon the  face of Athos, not even the respect he was accustomed to see on all faces.  Athos was dressed in black, with a simple lacing of silver. He wore the Holy  Ghost, the Garter, and the Golden Fleece, three orders of such importance, that  a king alone, or else a player, could wear them at once.

Mazarin rummaged a long time in his somewhat troubled memory to recall the name  he ought to give to this icy figure, but he did not succeed. <I am  told,> said he, at length, <you have a message from England for  me.>

And he sat down, dismissing Bernouin, who, in his quality of secretary, was  getting his pen ready.

<On the part of his majesty, the king of England, yes, your  eminence.>

<You speak very good French for an Englishman, monsieur,> said  Mazarin, graciously, looking through his fingers at the Holy Ghost, Garter, and  Golden Fleece, but more particularly at the face of the messenger.

<I am not an Englishman, but a Frenchman, monsieur le cardinal,>  replied Athos.

<It is remarkable that the king of England should choose a Frenchman for  his ambassador; it is an excellent augury. Your name, monsieur, if you  please.>

<Comte de la Fère,> replied Athos, bowing more slightly than the  ceremonial and pride of the all-powerful minister required.

Mazarin bent his shoulders, as if to say:—

<I do not know that name.>

Athos did not alter his carriage.

<And you come, monsieur,> continued Mazarin, <to tell  me\longdash>

<I come on the part of his majesty the king of Great Britain to announce  to the king of France>—Mazarin frowned—<to announce to  the king of France,> continued Athos, imperturbably, <the happy  restoration of his majesty Charles II to the throne of his ancestors.>

This shade did not escape his cunning eminence. Mazarin was too much accustomed  to mankind, not to see in the cold and almost haughty politeness of Athos, an  index of hostility, which was not of the temperature of that hot-house called a  court.

<You have powers, I suppose?> asked Mazarin, in a short, querulous  tone.

<Yes, monseigneur.> And the word <monseigneur> came so  painfully from the lips of Athos that it might be said it skinned them.

Athos took from an embroidered velvet bag which he carried under his doublet a  dispatch. The cardinal held out his hand for it. <Your pardon,  monseigneur,> said Athos. <My dispatch is for the king.>

<Since you are a Frenchman, monsieur, you ought to know the position of a  prime minister at the court of France.>

<There was a time,> replied Athos, <when I occupied myself  with the importance of prime ministers; but I have formed, long ago, a  resolution to treat no longer with any but the king.>

<Then, monsieur,> said Mazarin, who began to be irritated,  <you will neither see the minister nor the king.>

Mazarin rose. Athos replaced his dispatch in its bag, bowed gravely, and made  several steps towards the door. This coolness exasperated Mazarin. <What  strange diplomatic proceedings are these!> cried he. <Have we  returned to the times when Cromwell sent us bullies in the guise of charges  d'affaires? You want nothing, monsieur, but the steel cap on your head,  and a Bible at your girdle.>

<Monsieur,> said Athos, dryly, <I have never had, as you  have, the advantage of treating with Cromwell; and I have only seen his charges  d'affaires sword in hand; I am therefore ignorant of how he treated with  prime ministers. As for the king of England, Charles II., I know that when he  writes to his majesty King Louis XIV, he does not write to his eminence the  Cardinal Mazarin. I see no diplomacy in that distinction.>

<Ah!> cried Mazarin, raising his attenuated hand, and striking his  head, <I remember now!> Athos looked at him in astonishment.  <Yes, that is it!> said the cardinal, continuing to look at his  interlocutor; <yes, that is certainly it. I know you now, monsieur. Ah!  diavolo! I am no longer astonished.>

<In fact, I was astonished that, with your eminence's excellent  memory,> replied Athos, smiling, <you had not recognized me  before.>

<Always refractory and grumbling—monsieur—monsieur—What  do they call you? Stop—a name of a river—Potamos; no—the name  of an island—Naxos; no, per Giove!—the name of a  mountain—Athos! now I have it. Delighted to see you again, and to be no  longer at Rueil, where you and your damned companions made me pay ransom.  Fronde! still Fronde! accursed Fronde! Oh, what grudges! Why, monsieur, have  your antipathies survived mine? If any one has cause to complain, I think it  could not be you, who got out of the affair not only in a sound skin, but with  the cordon of the Holy Ghost around your neck.>

<My lord cardinal,> replied Athos, <permit me not to enter  into considerations of that kind. I have a mission to fulfill. Will you  facilitate the means of my fulfilling that mission, or will you not?>

<I am astonished,> said Mazarin,—quite delighted at having  recovered his memory, and bristling with malice,—<I am astonished,  Monsieur—Athos—that a Frondeur like you should have accepted a  mission for the Perfidious Mazarin, as used to be said in the good old  times\longdash> And Mazarin began to laugh, in spite of a painful cough,  which cut short his sentences, converting them into sobs.

<I have only accepted the mission near the king of France, monsieur le  cardinal,> retorted the comte, though with less asperity, for he thought  he had sufficiently the advantage to show himself moderate.

<And yet, Monsieur le Frondeur,> said Mazarin, gayly, <the  affair which you have taken in charge must, from the king\longdash>

<With which I have been given in charge, monseigneur. I do not run after  affairs.>

<Be it so. I say that this negotiation must pass through my hands. Let us  lose no precious time, then. Tell me the conditions.>

<I have had the honour of assuring your eminence that only the letter of  his majesty King Charles II contains the revelation of his wishes.>

<Pooh! you are ridiculous with your obstinacy, Monsieur Athos. It is  plain you have kept company with the Puritans yonder. As to your secret, I know  it better than you do; and you have done wrongly, perhaps, in not having shown  some respect for a very old and suffering man, who has laboured much during his  life, and kept the field for his ideas as bravely as you have for yours. You  will not communicate your letter to me? You will say nothing to me? Very well!  Come with me into my chamber; you shall speak to the king—and before the  king.—Now, then, one last word: who gave you the Fleece? I remember you  passed for having the Garter; but as to the Fleece, I do not know\longdash>

<Recently, my lord, Spain, on the occasion of the marriage of his majesty  Louis XIV, sent King Charles II a brevet of the Fleece in blank; Charles II immediately transmitted it to me, filling up the blank with my name.>

Mazarin arose, and leaning on the arm of Bernouin, he returned to his ruelle at  the moment the name of M. le Prince was being announced. The Prince de Condé,  the first prince of the blood, the conqueror of Rocroi, Lens, and Nordlingen,  was, in fact, entering the apartment of Monseigneur de Mazarin, followed by his  gentlemen, and had already saluted the king, when the prime minister raised his  curtain. Athos had time to see Raoul pressing the hand of the Comte de Guiche,  and send him a smile in return for his respectful bow. He had time, likewise,  to see the radiant countenance of the cardinal, when he perceived before him,  upon the table, an enormous heap of gold, which the Comte de Guiche had won in  a run of luck, after his eminence had confided his cards to him. So forgetting  ambassador, embassy and prince, his first thought was of the gold.  <What!> cried the old man—<all that—won?>

<Some fifty thousand crowns; yes, monseigneur,> replied the Comte  de Guiche, rising. <Must I give up my place to your eminence, or shall I  continue?>

<Give up! give up! you are mad. You would lose all you have won.  Peste!>

<My lord!> said the Prince de Condé, bowing.

<Good-evening, monsieur le prince,> said the minister, in a  careless tone; <it is very kind of you to visit an old sick  friend.>

<A friend!> murmured the Comte de la Fère, at witnessing with  stupor this monstrous alliance of words;—<friends! when the parties  are Condé and Mazarin!>

Mazarin seemed to divine the thoughts of the Frondeur, for he smiled upon him  with triumph, and immediately,—<Sire,> said he to the king,  <I have the honour of presenting to your majesty, Monsieur le Comte de la  Fère, ambassador from his Britannic majesty. An affair of state,  gentlemen,> added he, waving his hand to all who filled the chamber, and  who, the Prince de Condé at their head, all disappeared at the simple gesture.  Raoul, after a last look cast at the comte, followed M. de Condé. Philip of  Anjou and the queen appeared to be consulting about departing.

<A family affair,> said Mazarin, suddenly, detaining them in their  seats. <This gentleman is the bearer of a letter in which King Charles  II, completely restored to his throne, demands an alliance between Monsieur,  the brother of the king, and Mademoiselle Henrietta, grand-daughter of Henry  IV Will you remit your letter of credit to the king, monsieur le comte?>

Athos remained for a minute stupefied. How could the minister possibly know the  contents of the letter, which had never been out of his keeping for a single  instant? Nevertheless, always master of himself, he held out the dispatch to  the young king, Louis XIV, who took it with a blush. A solemn silence reigned  in the cardinal's chamber. It was only troubled by the dull sound of the  gold, which Mazarin, with his yellow, dry hand, piled up in a casket, whilst  the king was reading.  